% journal_results_main.tex — results at N = 20 on the nine analytical
% models (Sections 5.1-5.4). Numbers: results/journal/summary.json; tables
% generated by experiments/analyse_journal.py.

\section{Results}
\label{sec:results}
Table~\ref{tab:n20} gives the median nRMSE of every strategy on every model
at $N=20$, Table~\ref{tab:tests} the paired comparisons and
Fig.~\ref{fig:relative} the same relative to the Mat\'ern ARD GP with
confidence intervals.

\subsection{GP specification and variable transformation}
Estimating the kernel amplitude and one length scale per input (Mat\'ern ARD
versus the original kernel) reduced the nRMSE on six of the nine models, by
factors of 0.14 (physics) to 0.71 (piston), and did not change it on the
plume (both input models) and Ishigami models, where neither kernel can resolve the response
with 20 runs. On the physics model the original kernel, with its amplitude
fixed at the variance of the standardised output, cannot represent a
quadratic bowl whose values far exceed one standard deviation and reverts
towards the mean between runs; estimating the amplitude alone removes most
of this error, which the earlier analysis had attributed to the size of the
input domain (Section~\ref{sec:earlier}).

Fitting on the log scale, with log-transformed positive inputs, reduced the
nRMSE on all six models with positive outputs, by factors of 0.11
(borehole) to 0.75 (piston); all six reductions remain significant after
the Holm adjustment. The GP then interpolates a nearly additive function
of the log inputs, which is what products and power laws such as the
plume, borehole and OTL formulas become.

Adding a quadratic trend to the GP reduced the nRMSE on the physics model
by a factor of 400, to the level of the sparse PCE, and changed it by
less than 0.5\% on every other model: a GP with a trend loses practically
nothing when the trend is absent. On the log scale, the trend improved the
plume models (0.42 to 0.35 with uniform inputs, 0.37 to 0.28 with GUM-type
inputs), whose log concentration is close to quadratic in the inputs.

\subsection{Sparse PCE and committees}
Sparse PCE matched the trend GP on the physics model and the Mat\'ern GP on
the Ishigami model, and was 1.2 to 3.8 times worse than the GP on the other
seven models (1.1 to 12 times on the log scale); with 20 runs and up to
eight inputs, the corrected LOO error selects low-degree expansions that
cannot follow the borehole, OTL and flood responses as closely as the GP.

The CV-weighted committee was never better than the GP with a quadratic
trend on the same design. It tied on five models and was 1.7 to 2.4 times
worse on the other four: the borehole (both input models), OTL and flood
models. Its
advantage on the physics model in the earlier analysis is entirely that of
its degree-2 PCE member, which the trend GP reproduces (ratio 1.00). CV
weighting remained necessary for a committee: it was never worse than
uniform averaging and better on five models, by up to three orders of
magnitude on the physics model, where uniform averaging dilutes the exact
member with poor ones.

\subsection{Adaptive and targeted designs}
Space-filling adaptive designs gave no consistent gain at $N=20$. Adding
runs at the largest GP posterior SD never changed the nRMSE significantly
(ratios 0.86--1.22). The adaptive committee helped on the physics, borehole
and OTL models (0.61--0.80) and hurt on the plume model (1.25), where its
estimate of $P(C>\mathrm{IDLH})$ was off by a median of 0.17, against 0.02
for the passive committee. The spectral criterion UEigV was never better
than plain committee variance and was worse on three models (1.26--1.72).

Targeting the threshold changes the picture for exceedance probabilities.
AK-MCS reduced the median absolute error of $P(Y>T)$ significantly in 8 of
the 15 comparisons, by factors of 0.11 to 0.47, and was never significantly
worse. The largest gains were on the piston model (0.0099 to 0.0014 on the
log scale), the borehole model (0.0091 to 0.0009 on the raw scale) and the
GUM-type input models. Targeting comes at a cost elsewhere: on the
Ishigami model the AK-MCS design raised the median error of the propagated
mean from 8.5\% to 28\%, because its runs cluster near the threshold.

\subsection{Propagated quantities}
Table~\ref{tab:prop} gives the validation errors of the reported quantities
for the GP recommended for each model and for the committee. With the
recommended GP, the 95th percentile was within 0.2\% (OTL, flood) to 16\%
(Ishigami) of the reference, and $P(Y>T)$ within 0.0013 to 0.036 of its
reference value. The committee's 95th percentile was off by 6.9--31\% on
six of the nine models. At the IDLH threshold of the plume
model, the log-variable GP estimated $P(C>210\,\mathrm{mg\,m^{-3}})=0.123$
with a median absolute error of 0.009, against 0.023 for the raw-variable
GP.
